% !TeX root = ../教材题库_全部条目_转写工作稿.tex
% Source: XJ-B1-C04-EX-P102-Y0177
% 第4章 / 4.1 实数指数幂和幂函数 / 教材 p.102 / PDF 第109页
% OCR未校对：公式与文字均可能有误；未完成数学TeX转写；禁止直接发给学生
% 原题及相邻原解截图（核对用）：../题目原图/XJ-B1-C04-EX-P102-Y0177-01.webp
% 本文件只在转写完成后作为题干使用；OCR 条目尚不是可靠题干。
\begin{bankitem}{4.1 例题 例题8}{第4章 / 4.1 实数指数幂和幂函数 / 教材 p.102 / PDF 第109页}{XJ-B1-C04-EX-P102-Y0177}{OCR未校对：公式与文字均可能有误；未完成数学TeX转写；禁止直接发给学生}
\begin{ocrdraft}
[教材p.102 / PDF第109页]
\par
0河比较下列各组中两个数的大小：
\par
（1)1.51
\par
（2）1.50
\par
（3）1.5-1
\par
解
\par
（1）
\par
1.引
\par
，1.6凵可看作幂函数y一的两个函数值．
\par
该函数在[0，
\par
十。。）上递增，由于底数1.5<1.6，所以1.514
\par
（2）1.50
\par
，1.6“可看作幂函数3，一4的两个函数值．
\par
该函数在@，十亠）上递增，由于底数1.5<1.6，所以
\par
利用幂函数的单调性，
\par
1.50\ensuremath{\bullet}4<1.60口
\par
通过自变量的大小关系可
\par
（3）1.5-15，1.6一15可看作幂函数3，一的两个函数
\par
以判断相应函数值的大小
\par
值．该函数在（0，十。。）上递减，由于底数1.5<1.6，所以
\par
关系．
\par
1.5\ocrglyph{00AF}1\ensuremath{\bullet}5>1.6一1\ensuremath{\bullet}5
\end{ocrdraft}
\end{bankitem}
